\input amstex
\documentstyle{amsppt}
\magnification 1200
\pageno 35
\define\ogf{{G(F)}}
\define\ozf{Z_G(F)}
\define\oglf{{GL(2,E)}}
\define\oa{{A_G(F)}}
\define\oslf{{H(E)}}
\define\slf{{H(F)}}
\define\eps{\epsilon}
\define\ooo{{\Cal O}}
\define\ovo{\overline{\ooo}}
\define\off{{\Cal O}_F}
\define\offp{{\Cal O}_{F'}}
\define\qqq{\Bbb Q}
\define\ppp{\Bbb P}
\define\ccc{\Bbb C}
\define\rrr{\Bbb R}
\define\aaa{\Bbb A}
\define\nnn{\Bbb N}
\define\zzz{\Bbb Z}
\define\unt{\underline{T}}
\define\unx{\underline{x}}
\define\oga{{G(\aaa )}}
\define\oba{{B_G(\aaa )}}
\define\oza{{Z_G(\aaa )}}
\define\ogla{{GL(2,\aaa_L )}}
\define\osla{{H(\aaa_L )}}
\define\sla{H(\aaa )}
\define\oee{{\Cal O}_{E}}
\define\oeep{{\Cal O}_{E'}}
\define\VV{\Cal V}
\define\HH{\Cal H}
\hsize 6.5 true in
\vsize 8.7 true in
\vskip .2in
\subheading{\S 5 The Igusa expansion for $G_2$}
\vskip .1in
{\bf Notes on the derivation of the Igusa expansion}
\vskip .2in
Consider an $F$-morphism 
$$
\phi:Y\rightarrow {\Bbb A},
$$
smooth outside the ``special fiber'' $\phi^{-1}(0)$,
from a smooth $F$-variety $Y$ to the affine line over $F$.
When we regard $Y(F)$ and ${\Bbb A}(F)=F$ as $F$-manifolds
in the p-adic topology it will be explicitly stated.
The ``default'' topology (to use the computer terminology)
on $Y$ is the Zariski topology. Let $f\in C_c^\infty(Y(F))$,
considering $Y(F)$ as a p-adic manifold, and consider
$$
F(\lambda):=
\int_{\phi^{-1}(\lambda)(F)}f(y)|\omega_\lambda|,
$$
for $\lambda\in F-\{0\}$. If $f$ is locally a product of 
quasi-characters on a suitable local coordinate system of
a neighborhoood of $\phi^{-1}(0)$ in the manifold $Y(F)$ then Igusa 
has given an ``asymptotic expansion'' for $F(\lambda)$
as $\lambda\rightarrow 0$. Actually, since the topology on 
$Y(F)$ is totally disconnected, this ``asymptotic expansion''
is actually an equality valid for all $\lambda$ in a neighborhood
of $0$. We explain this next.
\par
Assume that $Y(F)$ may be covered by patched
$U$ on which there exist coordinates $\mu_1,...,\mu_n$
such that 
\newline
(5.1) $\phi^{-1}(0)$ has normal crossings; more specifically,
the irreducible components $E_i$ of $\phi^{-1}(0)\cap U$ 
are of the form $\{\mu_i =0\}$ on $U$ and
$$
\lambda=\alpha\prod_i\mu_i^{a_i},
$$
for some integral $a_i>0$, where $\alpha$ is a non-vanishing 
function in a p-adic neighborhood of $0$ and 
$|\alpha|=q^{-m(\alpha)}$ is constant for all sufficiently
small $\lambda$.
\newline
(5.2) On $U$, $f(y)$ is of the form
$$
f(\mu_1,...,\mu_n)=\gamma \prod_i\kappa_i(\mu_i),
$$
where $\gamma$ is constant for sufficiently small $\lambda$.
\newline
(5.3) the form $\omega$ is given on $U$ by
$$
\omega=\prod_i\mu_i^{b_i-1}d\mu_1...d\mu_n.
$$
(5.4) We {\bf assume} that $f$ vanishes off $U$ and that if $b_i=0$ then
$\kappa_i\not= 1$, $1\leq i\leq n$.
\par
Let $F(\lambda)$ be given by (0.2) above. 
We will sketch the derivation of an equality of the form
$$
F(\lambda) =\sum_{\theta,\beta}\theta(\lambda)|\lambda|^{\beta -1}
\sum_{r=1}^{e(\theta,\beta)}(-\log_q|\lambda|)^{r-
1}F_r(\theta,\beta,f).
\tag 5.5
$$
The following lemma helps us to compute $F_1$:
\proclaim{Lemma 5.6} Let
$$
G(\lambda) =\sum_{\theta,\beta}\sum_{r=1}^{e(\theta,\beta)}
a(\theta,\beta,r)(-\log_q|\lambda|)^{r-1}\theta(\lambda)|\lambda|^{\beta -1},
$$
where the $a(\theta,\beta,r)$ are constants. If 
$a(\theta,\beta,r)=0$ for $r>1$ then
the coefficients $b_n$ in the power series of the 
truncated Mellin transform
$$
M_\epsilon G(s):=\int_{|\lambda|\leq \epsilon}
G(\lambda)|\lambda|^sd\lambda =\sum_{n\geq 0}b_nq^{-ns},
\ \ \ \ \ \epsilon=q^{-M},
$$
are given by
$$
b_n=
\cases
vol({\Cal O}_F^\times,d\lambda)
\sum_{\Sb \theta,\beta\\ \theta\text{ unramified}\endSb}
a(\theta,\beta,1)(\theta(\pi)q^{-\beta})^n,&n\geq M,\\
0,& n<M.
\endcases
$$
\endproclaim
The proof of this is left to the reader.
\par
Now, to verify the formula (5.5) for $F(\lambda)$ we may 
verify it for its truncated Mellin transform
and use the lemma above to check the coefficients of
both sides of the resulting power series in
$q^{-s}$. The formula we have in mind for $F_1$ is stated below.
First, we need some notation: in the notation of the assumptions
(5.1-4) above, let $c_j=c_j(a_1,...,a_n)$ denote
the coefficients for which the principal part of
$\prod_i{1\over 1-t^{a_i}}$ at $t=1$ is
$$
\sum_j{c_j\over (1-t)^j}.
$$
Let
$$
A(x)=\sum_j {c_j\over (j-1)!}(x+1)...(x+j-1),
$$
and define $A_i(x)$ by $A(x-y)=\sum_j x^{j-1}A_j(y)$,
so $A_1(y)$ is the ``constant term'' of this polynomial in $x$.
(The function $A(x)$ arises via the undergraduate calculus equality
$$
{(j-1)!\over (1-t)^j}
=\sum_{n\geq 0}(n+1)(n+2)...(n+j-1)t^n.)
$$
\proclaim{Lemma 5.7} Suppose that $Y(F)$, $f$, $\omega$, and 
$\lambda$ satisfy the conditions of (5.1-4). Suppose furthermore 
that the support of $f$ is contained in one such patch
$U$, where
$$
U=\{(\mu_1,...,\mu_n)\in F^n\ |\ |\mu_i|\leq q^{-M_i}\}.
$$
Finally, suppose that we reindexed the $\mu_i$, if 
necessary, so that 
$$
\theta^{a_i}=\kappa_i,\ \ \ \ \beta=b_i/a_i,\ \ \ \ 
1\leq i\leq p,
$$
and for each $i>p$ at least one of these equalities fails to hold.
(This $p$ is of course different from the fixed prime $p$.)
Then $F(\lambda)$ satisfies the formula (5.5)
where $F_1(\theta,\beta,f)$ is given by the value of
$$
\int_{D_p\cap U}A_1(M)h_p\prod_{i=p+1}^n|\mu_i|^{z_i}|\nu_p|,
$$
at $z_{p+1}=...=z_n=0$, where
$$
M=m(\alpha)+\sum_{i\leq p}a_iM_i +\sum_{i>p}a_i(-\log_q|\mu_i|),
$$
$$
h_p={\gamma\over \theta(\alpha)}
\prod_{i=p+1}^n\kappa_i(\mu_i)\theta(\mu_i)^{-a_i},
$$
and 
$$
\nu_p={1\over |\alpha|^\beta}
\prod_{i=p+1}^n|\mu_i|^{b_i-\beta a_i-1}d\mu_{p+1}...d\mu_n.
$$
Moreover, if, in addition, any one of the characters $\kappa_i\theta^{-a_i}$
is ramified then $F_1(\theta,\beta,f)=0$.
\endproclaim
We now sketch the proof of this lemma following \cite{L}.
By our assumptions,
$$
\split
M_\epsilon (F\theta^{-1})(s)
&={|\alpha|^s\over \theta(\alpha)}
\int_Uf(y)\prod_i\theta(\mu_i)^{-a_i}|\mu_i|^{a_is+b_i-1}
d\mu_1...d\mu_n\\
&={|\alpha|^s\over \theta(\alpha)}\gamma
\prod_i\int_{|\mu|\leq q^{-M_i}}
\kappa_i(\mu)\theta(\mu)^{-a_i}|\mu|^{a_is+b_i-1}d\mu,
\endsplit
\tag 5.8
$$
since $|\omega_\lambda ||d\lambda |
=\prod_{i=1}^n|\mu_i|^{b_i-1}d\mu_1...d\mu_n$.
Recall that 
$$
\int_{|\mu|\leq q^{-M_i}}
\kappa'(\mu)|\mu|^zd\mu =0
$$
if $\kappa'$ is a ramified character. Therefore we may assume 
from now on that every character $\kappa_i(\mu)\theta(\mu)^{-a_i}$
is unramified and we let
$$
\epsilon_i:=\kappa_i(\pi)\theta(\pi)^{-a_i},\ \ \ \ 
1\leq i\leq n,
$$
where $\pi$ denotes a fixed generator of the maximal ideal
of the local ring of integers ${\Cal O}_F$. We may write
$$
M_\epsilon (F\theta^{-1})(s)
=\gamma{|\alpha|^s\over \theta(\alpha)}
(1-q^{-1})^n\prod_i
{(\epsilon_iq^{-a_is-b_i})^{M_i}\over 
1-\epsilon_iq^{-a_is-b_i}}.
$$
Recall that we have reindexed the $\mu_i$ so that 
$$
\theta^{a_i}=\kappa_i,\ \ \ \ \beta=b_i/a_i,\ \ \ \ 
1\leq i\leq p,
$$
and for each $i>p$ at least one of these equalities fails to hold.
Then
$$
M_\epsilon (F\theta^{-1})(s)
=M(s,z_{p+1},...,z_n),
$$
at $z_{p+1}=...=z_n=0$, where
$$
\split
&M(s,z_{p+1},...,z_n)\\
&\quad\quad =(1-q^{-1})^p
[\prod_{1\leq i\leq p}
{(q^{-a_is-b_i})^{M_i}\over 
1-q^{-a_is-b_i}}]|\alpha|^{s+\beta}
\int_{\Sb |\mu_i|\leq q^{-M_j}\\ p+1\leq i\leq n\endSb}
\times\\
&\quad\quad\quad\quad \times
h_p|\alpha|^{-\beta}
\prod_{p+1\leq i\leq n}|\mu_i|^{z_i+a_is+b_i-1}
d\mu_{p+1}...d\mu_n\\
&\quad\quad
=(1-q^{-1})^p|\alpha|^{s+\beta}{\gamma\over \theta(\alpha)}
[\prod_{1\leq i\leq p}
{(q^{-a_is-b_i})^{M_i}\over 
1-q^{-a_is-b_i}}]\times \\
&\quad\quad \quad\quad \times
|\alpha|^{-\beta}\int_{\Sb |\mu_i|\leq q^{-M_j}\\ p+1\leq j\leq n\endSb}
\prod_{p+1\leq i\leq n}|\mu_i|^{a_i(s+\beta)}
\kappa_i(\mu_i)\theta(\mu_i)^{-a_i}|\mu_i|^{z_i-a_i\beta+b_i-1}
d\mu_{p+1}...d\mu_n\\
&\quad\quad
=(1-q^{-1})^p|\alpha|^{s+\beta}[\prod_{1\leq i\leq p}
{(q^{-a_is-b_i})^{M_i}\over 
1-q^{-a_is-b_i}}]\times \\
&\quad\quad \quad\quad \times
\int_{\Sb |\mu_i|\leq q^{-M_j}\\ p+1\leq i\leq n\endSb}
\prod_{p+1\leq i\leq n}|\mu_i|^{z_i+a_i(s+\beta)}h_p|\nu_p|.
\endsplit
\tag 5.9
$$
The last integral may be interpreted as a principal-value
integral of a function on $Y(F)$ over a variety of the form
$D_p=\cap_{1\leq i\leq p}E_i$, where $E_i$ is an irreducible 
component of $\phi^{-1}(0)$ which satisfies
$E_i\cap U=\{(\mu_1,...,\mu_n)\ |\ \mu_i=0\}$. Moreover, this 
multiple integral is a rational function of $q^{-s}$.
\par
To see what this rational function is, let for the moment
$t=q^{-s-\beta}$ and recall that we have defined the $c_j$ in 
such a way that the principal part of
$$
\prod_{1\leq i\leq p}
{(q^{-a_i(s+\beta)})^{M_i}\over 
1-q^{-a_i(s+\beta)}}=
\prod_{1\leq i\leq p}
{t^{a_iM_i}\over 
1-t^{a_i}}
$$
at $t=1$ is
$$
t^{\sum_{1\leq i\leq p}a_iM_i}
\sum_{1\leq i\leq p}{c_j\over (j-1)!}(1-t)^{-j}
=\prod_{p+1\leq i\leq n}|\mu_i|^{a_i(s+\beta)}|\alpha|^{-s-\beta}
t^M\sum_{n\geq 0}A(n)t^{n},
$$
where $M$ is as in the statement of the lemma.
(Here we have multiplied and divided by
$\prod_{p+1\leq i\leq n}|\mu_i|^{a_i(s+\beta)}
|\alpha|^{-s-\beta}$ to account for all of $M$.)
The coefficient of $t^n$ in the latter expression is,
for $n\geq M$,
$$
A(n-M)\prod_{p+1\leq i\leq n}|\mu_i|^{a_i(s+\beta)}|\alpha|^{-s-\beta}
=\prod_{p+1\leq i\leq n}|\mu_i|^{a_i(s+\beta)}|\alpha|^{-s-\beta}
\sum_rn^{r-1}A_r(M),
$$
and, for $n<M$, the coefficient of $t^n$ is $0$. From this the 
lemma follows. 
\vskip .2in
{\bf Geometry of the Igusa expansion}
\vskip .1in
Let $E$ be an irreducible divisor in $Y''\subset G\times S''$, 
which we know is regular in 
codimension one, for which the image of $E$ under $\pi$ is contained in the
Zariski closure of a unipotent orbit ${\Cal O}$ in $G$. Such a divisor is 
called an {\bf ${\Cal O}$-divisor} if ${\Cal O}$ is minimal (with respect to 
the partial ordering given by inclusion of Zariski closures) with this 
property. It will be called an {\bf ${\Cal O}(F)$-divisor} if $E$
is defined over $F$ and
$\pi(E(F))$ meets ${\Cal O}(F)$ in a dense subset of $\ooo(F)$ 
as p-adic manifolds.
Let $E$ be an ${\Cal O}(F)$-divisor.
\par
First, we record an observation about
the fibration $\pi:E\rightarrow \overline{\ooo}$.
As a corollary to \cite{DG, ch II, Prop. 3.1}, we have
\proclaim{Lemma 5.10} 
(a) The morphism $\pi:E\cap Z_S \rightarrow \overline{\ooo}$ is 
smooth of relative 
dimension $2r$, provided $\ooo$ is $r$-regular. 

(b) $\pi(E\cap Z_S)\subset \ooo$.


(c) Let 
$$
\{1\}\subset \overline{\Cal O}_1={\Cal O}_1\cup \{1\}
\subset ...\subset \overline{\Cal O}={\Cal O}\cup {\Cal O}_m,
$$
where each ${\Cal O}_i$ is a unipotent orbit of $G$,
and let $E_i=\pi^{-1}({\Cal O}_i)\cap E$. Then $\pi|_{E}:E_i
\rightarrow {\Cal O}_i$ is faithfully flat. 
\endproclaim
We may apply Igusa's expansion to
$$
F(\lambda)=
\int_{\phi^{-1}(\lambda)(F)}f(y)|\omega_\lambda|,
\tag 5.11
$$
where $f$ is any function on $Y_C(F)$ for which (a) $f$ is compactly
supported on the p-adic manifold $Y_C(F)$ and (b) $\phi$, $|\omega_\lambda|$,
$Y_C$, and $f$ satisfy the conditions of Igusa data in \cite{LS0}. 
Since $f$ is compactly supported, so is $f|_{\ooo}$. On each fiber 
$\tilde{E}(u)$ the pull-back $\pi^*(f)$ takes a constant value $f(u)$.
By \cite{H1, Ch I}, the regular function $\kappa(t_*)$ extends
from $E(u)$ to a regular function on $\tilde{E}(u)$. By the lemma
above, we have
$$
\int_{\overline{\ooo}(F)}f \cdot
PV\int_{\tilde{E}(u)(F)}\kappa(\tilde{t}_*)
|\omega_{\tilde{E}(u)}||\omega_{\ovo}|=
PV\int_{\tilde{E}(F)}\tilde{\pi}^*(f) \kappa(\tilde{t}_*)
|\omega_{\tilde{E}}|,
\tag 5.12
$$
provided everything in sight satisfies the conditions (5.1-4).
Recall that, by definition, there are disjoint patches 
$U_i\subset \tilde{Y}_1(F)$,
each analytically isomorphic to a finite box in affine space, such that
$$
PV\int_{\tilde{E}(F)}\tilde{\pi}^*(f) \kappa(\tilde{t}_*)
|\omega_{\tilde{E}}|=
\sum_iPV\int_{\tilde{E}(F)\cap U_i}\tilde{\pi}^*(f) \kappa(\tilde{t}_*)
|\omega_{\tilde{E}}|,
\tag 5.13
$$
provided that the conditions (5.1-4) hold.
(By the way, the computations in \S 4 show that in all
the cases looked at in this paper,
$b_{{E}(F)}>0$. In particular, $b_{{E}(F)}\not= 0$.)

\proclaim{Lemma 5.14}
Let $G=G_2$.
Let $E$ be a base ${\Cal O}(F)$-divisor of the 
variety $Y_C=\tilde{Y}_1$. If $E$ is $F$-rational, $E(F)\not= \emptyset$,
then $e(\theta_E,\beta_E)=1$. In particular, spurious divisors
do not contribute to the Igusa expansion.
\endproclaim
\demo{proof} The proof is case-by-case using local coordinates in each case.
Assume $e(\theta_E,\beta_E)\geq 2$. By \cite{L}, the
principal value integral contributing to 
$F_1(\theta_E,\beta_E,f_\kappa)$ is
$$
\int_{\ooo(F)}f\int_{D_u(F)}\kappa(\tilde{t}_*)
|\omega_{D_u}||\omega_{\ooo}|,
$$
where $D=\cap_{i=1}^{e(\theta_E,\beta_E)}E_i$, 
where $\theta_E=\theta_i$ and $\beta_E=\beta_i$,
and $D$ is $F$-rational (though the $E_i$ themselves need not be
$F$-rational). 

Claim: The homogeneity degree of $|\omega_{D_u}|$, as a function
of the tangent direction $X$, is $dim(D)-dim(\ooo)<dim(E)-dim(\ooo)=2r$.

To prove this, we look at each divisor in the inverse image of
$\{\lambda =0\}$. 

We may omit the regular and 1-regular cases since they were already
dealt with in \cite{H1}, \cite{LS2}.

Let ${\Cal O}={\Cal O}_2$ denote a 2-regular orbit and let 
$E=E_{\emptyset, \alpha,\delta}$. 
First suppose $e(\theta_E,\beta_E)=2$.
From \S 7 below, the form on the $F$-rational fiber
$E_{\emptyset,\alpha,\delta}(u)$, $u\in {\Cal O}$, is given 
on $Y''_{\emptyset,\alpha,\delta}$ by
$$
\omega_{E_{\emptyset,\alpha,\delta}(u)}
=gdz(\beta)dx(\zeta)dv(\alpha)dv(\zeta),
\tag 5.15
$$
where $g$ is a regular function.
(Though these coordinates may not be $F$-rational, we
may ``average'' them over the image of the Galois group in $\Omega$
using the calculations of \S 2 and then we make a change
of variables in the form to get a form in $F$-rational coordinates
satisfying (5.1-4).)
By (1.22), this form is homogeneous
of degree 2 as a function of $X$. Let $E'$ denote another 
$F$-rational ${\Cal O}_2$-rational divisor with
$\theta_E=\theta_{E'}$, $\beta_E =\beta_{E'}$. By Hironaka's theorem,
$E$ intersects $E'$ normally in $Y_C$. The intersection $D=E\cap E'$
maps onto ${\Cal O}$ under $\pi$ by \cite{DG, chap II, Prop 3.1}. 
The variety $D$ must have codimension 1 in $E$ and will have
an Igusa form on 
a fiber $D(u)$ which satisfies
$$
\omega_{D(u)}=g_Ddx dv_1dv_2,
$$
where $x$ is a coordinate on $Z$, $v_1,v_2$ are coordinates on the
$N_+$ components, and $g_D$ is regular. Therefore, $\omega_{D(u)}$
is homogeneous of degree $<2$ as a function of $X$. 

Now suppose $e(\theta_E,\beta_E)=3$. The argument is similar,
hence omitted, in this case. The claim follows in the
2-regular case.

Let ${\Cal O}={\Cal O}_3$ denote a 3-regular orbit and let 
$E=E_{\emptyset, \beta,\zeta}$. 
First suppose $e(\theta_E,\beta_E)=2$.
From \S 8, the form on the $F$-rational fiber
$E_{\emptyset,\beta,\zeta}(u)$, $u\in {\Cal O}$, is given 
on $Y''_{\emptyset,\beta,\zeta}$ by
$$
\omega_{E_{\emptyset,\beta,\zeta}(u)}
=h dw(\epsilon)dw(\zeta)dx(\zeta)dv(\beta)dv(\epsilon)dv(\zeta),
\tag 5.16
$$
where $h$ is a regular function.
(Again, these coordinates may be ``averaged to obtain $F$-rational
coordinates as in the 2-regular case.)
By (1.22), this form is homogeneous
of degree 3 as a function of $X$. Let $E'$ denote another 
$F$-rational ${\Cal O}_3$-rational divisor with
$\theta_E=\theta_{E'}$, $\beta_E =\beta_{E'}$. By Hironaka's theorem,
$E$ intersects $E'$ normally in $Y_C$. The intersection $D=E\cap E'$
maps onto ${\Cal O}$ under $\pi$ by \cite{DG, chap II, Prop 3.1}. 
The variety $D$ must have codimension 1 in $E$ and will have Igusa form on 
a fiber $D(u)$ which satisfies
$$
\omega_{D(u)}=h_Ddx dy dv_1dv_2dv_3,
$$
where $x,y$ are coordinates on $Z$, $v_1,v_2,v_3$ are coordinates on the
$N_+$ components, and $h_D$ is regular. Therefore, $\omega_{D(u)}$
is homogeneous of degree $<3$ as a function of $X$. 

Now suppose $e(\theta_E,\beta_E)=3$. The argument is similar,
hence omitted, in this case. The claim follows in the
3-regular case.
This proves the claim for all the orbits. 

However, the claim 
contradicts the degree of homogeneity of the $\ooo$-germ. \qed
\enddemo
\proclaim{Lemma 5.17} Let $G=G_2$.
Let $E$ be a base ${\Cal O}(F)$-divisor of the 
variety $Y_1$, $\tilde{E}$ its preimage under 
$Y_C\rightarrow Y_1$.
In this case $F_1(\theta_E,\beta_E,f_\kappa)$ is given 
(up to a constant) by 
$$
PV\int_{\tilde{E}(F)}\tilde{\pi}^*(f_\kappa) |\omega_{\tilde{E}}|
=\int_{\ovo(F)_{st}}f
\int_{\tilde{E}_u(F)}\kappa(\tilde{t}_*)
|\omega_{\tilde{E}_u}||\omega_{\ovo}|,
$$
where $\omega_{\ovo_{st}}$ 
is a sum of forms of maximal degree on the 
equidimensional pieces ${\Cal O}_i$ in Lemma 5.10 of the stable
orbit $\ovo(F)_{st}=\cup_i{\Cal O}_i(F)$.
\endproclaim
\demo{proof} Because of the previous lemma, (0.1), and formula 0.6,
it suffices to check that the conditions of \cite{LS0} are met. This
has essentially already been done in \cite{L},
\cite{H1, ch I, proof of Prop. 5.2},
but we fill in some details regarding the function $\kappa(\tilde{t}_*)$ 
and the cocycle-valued function $t_*$.
\par
Some information about $t_*$ is given following formula 0.6 in \S 0.
More explicit information about
$t_\sigma :Y_1(B_+,B_-)\rightarrow T_G(\overline{F})$, where
$Y_1(B_+,B_-)$ is defined in \S 1, is given in
\cite{H1, ch I, Corollary 5.4}. Since the
only $F$-form of $G_2$ over a p-adic field is the split form \cite{T},
we have $G=G_2=G_{in}=G_{qs}$, $T_G=T_{in}, \ \underline{T}=T_{qs}$, and
$\underline{B}=B_{qs}$, the notation being as in \cite{H1}. Let 
$\alpha\in \Delta$ denote a simple root,
If $\sigma=\sigma_\alpha\in \Gamma$ acts on
$T_G(\overline{F})$ by conjugation by the
reflection $s_\alpha\in \Omega$, then
we choose $g$, $h$ as in 
\cite{H1, ch I, \S 5}. Taking
$A_\sigma=w_\sigma^{-1}=s_\alpha \in N_G(T_G)$.
in \cite{H1, ch I, Corollary 5.7}, we find there is a ``constant''
$t_1\in T_G(\overline{F})$ independent of the star such that at a point
$$
(t^g,((B(W))_{W\in{\Cal W}},(z(\alpha))_{\alpha\in\Delta},n)
\in\chi^{-1}(C)\times S_1(B_+,B_-)\times N_+,
$$
the value $t_\sigma$ is determined by
$$
s_\alpha \epsilon_{-\alpha}(-\xi)n^{-1}w_\sigma^{-1}s_\alpha^{-1}\in 
N_+^- N_+ (t_1t_\sigma)^{w_\sigma^{-1}s_\alpha^{-1}}.
\tag 5.18
$$
Take the transpose of both sides of Equation 5.5 in
\cite{H1, ch I}. The resulting equation,
along with (5.10) implies that 
$$
t_1t_\sigma = \alpha^\vee(\xi),
$$
where $\alpha^{\vee}$ is its coroot, and $\xi$ satisfies (5.2)
as a function of the coordinates of $Y''(B_+,B_-)\times N_+$. Choose
$B_-, B_+$ such that these coordinates are $F$-rational.
From the appendix, we know that there is a natural bijection
$$
H^1(\Gamma,T_G)\rightarrow \zzz/ m\zzz,
$$
where $m\in \{2,3\}$. Since each such $H^1(\Gamma,T_G)$ is cyclic,
$\kappa(\tilde{t}_*)$ is determined by its value on a generator. This,
in turn, determines a character of $\zzz/m\zzz$ and
hence by local class field theory a
character $\theta$ of $F^\times/NE_{T_G}^\times$ (where $E_{T_G}$
denotes the Galois closure of the splitting field of $T_G$. One may
construct a (non-canonical) map
$\phi:H^1(\Gamma,T_G) \rightarrow F^\times$
such that $\kappa(\tilde{t}_*)=\theta(\phi(\tilde{t}_*))$
is a monomial in the coordinates of $Y''(B_+,B_-)\times N_+$.
This allows us to conclude that the
conditions of (5.1-4) have been met.  \qed
\enddemo

We close this section by recalling an important reduction lemma
due to T. Hales. First, observe that for each $r$-regular
orbit ${\Cal O}$, $r>0$,
every function $f\in C_c^\infty(G(F))$ may be written in the form $f=f_1+f_2$
with $f_1$ supported off the Zariski closure of ${\Cal O}$ in $G$
and $f_2$ whose support does not meet the orbits outside the 
Zariski closure of ${\Cal O}$ in $G$.
\proclaim{Reduction Lemma 5.19 (Hales)} Suppose that for each $(T,\kappa)$ 
associated to $H$ 
the germs $\Gamma_{{\Cal O}}^{(T,\kappa)}(\gamma)$ have the form
$$
\Gamma_{{\Cal O}}^{(T,\kappa)}(exp(\lambda^2 X))=
|\lambda|^{2r}m_{{\Cal O}}^{(T,\kappa)}(X),
$$
for $\lambda$ small, where $m_{{\Cal O}}^{(T,\kappa)}(X)$ satisfies
$$
m_{{\Cal O}}^{(T,\kappa)}(X)=m_{{\Cal O}}^{(T,\kappa)}(X+Y),
$$
for $X,Y\in {\frak t}_G$ sufficiently small, $X$ regular,
$X+Y$ regular, and $Y$ belonging to the Lie algebra of the center
of an endoscopic group $H\subset G$ satisfying $Cent(Y,G)=H$. 
Then $(G,H)$ admits local
$\Delta$-transfer at the identity for all functions $f\in C_c^\infty(G(F))$
whose support does not meet the orbits outside the 
Zariski closure of ${\Cal O}$ in $G$. 
\endproclaim
\demo{proof} Follows \cite{H3, Lemma 6.1}. (See also \cite{H2, Proposition 
2.11}.)  \qed
\enddemo

This lemma implicitly uses \cite{LS2}.
\enddocument
